\subsection{模的支撑集}

定义 5. 设 A 是一个环，M 是一个 A-模。A 的满足 \(M_{p} \neq 0\) 的素理想 p 组成之集称为 M 的支撑集，记作 Supp(M)。

由于 A 的每个极大理想都是素理想，由 §3 第 3 节定理 1 的推论 2 立即得出：为使 A-模 M 等于 0，必须且只需 \(\operatorname{Supp}(\mathbf{M}) = \varnothing\)。

例. 设 a 是 A 的一个理想；用第 3 节的记号，我们有

\[
\mathbf {V} (\mathfrak {a}) = \operatorname{Supp} (\mathrm{A} / \mathfrak {a}).\tag{9}
\]

若 \(p\) 是 \(A\) 的满足 \(a \not\subset p\) 的素理想，则 \((A/a)_p = 0\)（\(\S 3\) 第 1 节注 3）；反之，若 \(a \subset p\)，则 \(aA_p\) 含于 \(pA_p\)（\(A_p\) 的极大理想），且 \((A/a)_p\) 同构于 \(A_p/aA_p\)，从而非零（\(\S 3\) 第 1 节命题 3）；由此得我们的断言。

特别地，\(\operatorname{Supp}(\mathbf{A}) = \operatorname{Spec}(\mathbf{A})\)。

命题 16. 设 A 是一个环，M 是一个 A-模。

\begin{enumerate}
\def\labelenumi{(\roman{enumi})}
\tightlist
\item
  若 \(\mathbf{N}\) 是 \(\mathbf{M}\) 的一个子模，则
\end{enumerate}

\[
\operatorname{Supp} (\mathbf {M}) = \operatorname{Supp} (\mathbf {N}) \cup \operatorname{Supp} (\mathbf {M} / \mathbf {N}).
\]

\begin{enumerate}
\def\labelenumi{(\roman{enumi})}
\setcounter{enumi}{1}
\tightlist
\item
  若 \(\mathbf{M}\) 是子模族 \((\mathbf{N}_{\iota})_{\iota \in \mathbf{I}}\) 之和，则
\end{enumerate}

\[
\operatorname{Supp} (\mathbf {M}) = \bigcup_ {i \in I} \operatorname{Supp} \left(\mathrm{N} _ {i}\right).
\]

(i) 由正合列 \(0 \to \mathbf{N} \to \mathbf{M} \to \mathbf{M} / \mathbf{N} \to 0\)，对每个素理想 \(\mathfrak{p}\)（\(\mathbf{A}\) 的）我们导出正合列

\[
0 \rightarrow \mathrm{N} _ {p} \rightarrow \mathrm{M} _ {p} \rightarrow (\mathrm{M} / \mathrm{N}) _ {p} \rightarrow 0
\]

（§2 第 4 节定理 1）。为使 \(\mathbf{M}_{\mathfrak{p}}\) 化为 0，必须且只需 \(\mathrm{N}_{\mathfrak{p}}\) 与 \((\mathrm{M} / \mathrm{N})_{\mathfrak{p}}\) 都化为 0。换言之，关系 \(\mathfrak{p} \notin \operatorname{Supp}(\mathbf{M})\) 等价于`` \(\mathfrak{p} \notin \operatorname{Supp}(\mathbf{N})\) 且 \(\mathfrak{p} \notin \operatorname{Supp}(\mathbf{M} / \mathbf{N})\) ''，这就证明了 (i)。

(ii) 对 A 的每个素理想 p，\(M_{p}\) 是子模族 \((\mathbf{N}_{\iota})_{\mathfrak{p}}\) 之和（§2 第 4 节）。说 \(M_{p} \neq 0\) 成立，意味着存在 \(\iota \in I\) 使 \((\mathbf{N}_{\iota})_{\mathfrak{p}} \neq 0\)，由此得 (ii)。

推论. 设 A 是一个环，M 是一个 A-模，\((m_{i})_{i\in I}\) 是 M 的一个生成系，\(a_{i}\) 是 \(m_{i}\) 的零化子。则 Supp(M) = \(\bigcup_{i\in I} V(a_{i})\)。

由命题 16 (ii)，\(\operatorname{Supp}(\mathbf{M}) = \bigcup_{i \in I} \operatorname{Supp}(Am_i)\)。另一方面，\(Am_i\) 同构于 A-模 \(A/a_i\)，且我们已见

\[
\operatorname{Supp} (\mathrm{A} / \mathfrak {a} _ {\iota}) = \mathrm{V} (\mathfrak {a} _ {\iota})
\]

（上面的例）。

命题 17. 设 A 是一个环，M 是一个 A-模，a 是它的零化子；若 M 是有限生成的，则 \(\operatorname{Supp}(\mathbf{M}) = \mathrm{V}(\mathfrak{a})\)，因而 \(\operatorname{Supp}(\mathbf{M})\) 在 \(\operatorname{Spec}(\mathbf{A})\) 中是闭的。设 \((m_i)_{1 \leqslant i \leqslant n}\) 是 M 的一个生成系，\(\mathfrak{a}_i\) 是 \(m_i\) 的零化子；则 \(\mathfrak{a} = \bigcap_{i=1}^{n} \mathfrak{a}_i\)，从而 \(\mathrm{V}(\mathfrak{a}) = \bigcup_{i=1}^{n} \mathrm{V}(\mathfrak{a}_i)\)（第 3 节公式 (3)），该命题由命题 16 的推论得出。

推论 1. 设 A 是一个环，M 是一个有限生成 A-模，a 是 A 的一个元素。为使 a 属于 M 的支撑集的每个素理想，必须且只需 M 上比率为 a 的位似是幂零的。

由命题 17，属于 Supp(M) 的素理想之交是 M 的零化子 a 的根（第 3 节命题 11 (ii)）。说 a 属于该根，等价于说存在一个幂 \(a^{k} \in a\)，从而 \(a^{k}M = 0\)。

推论 2. 设 A 是 Noether 环，M 是一个有限生成 A-模，a 是 A 的一个理想。为使 \(\operatorname{Supp}(\mathbf{M}) \subset \mathrm{V}(\mathfrak{a})\)，必须且只需存在整数 \(k\) 使 \(\mathfrak{a}^k\mathbf{M} = 0\)。

若 b 是 M 的零化子，则由命题 17，关系 \(\operatorname{Supp}(\mathbf{M}) \subset \mathbf{V}(\mathfrak{a})\) 等价于 \(\mathbf{V}(\mathfrak{b}) \subset \mathbf{V}(\mathfrak{a})\)，从而等价于 \(\mathfrak{a} \subset \mathfrak{r}(\mathfrak{b})\)，其中 \(\mathfrak{r}(\mathfrak{b})\) 是 b 的根（第 3 节命题 11 的推论 2）。由于 A 是 Noether 的，该条件也等价于存在整数 \(k > 0\) 使 \(a^{k} \subset b\)（§2 第 6 节命题 15）。

命题 18. 设 M、M' 是环 A 上的两个有限生成模；则

\[
\operatorname{Supp} (\mathbf {M} \otimes_ {\mathbf {A}} \mathbf {M} ^ {\prime}) = \operatorname{Supp} (\mathbf {M}) \cap \operatorname{Supp} (\mathbf {M} ^ {\prime}).\tag{10}
\]

需要证明：若 p 是 A 的素理想，则关系 \((\mathbf{M} \otimes_{\mathbf{A}} \mathbf{M}')_{p} \neq 0\) 与`` \(M_{p} \neq 0\) 且 \(M'_{p} \neq 0\) ''等价。由于 \(A_{p}\) -模 \(M_{p} \otimes_{A_{p}} M'_{p}\) 与 \((\mathbf{M} \otimes_{\mathbf{A}} \mathbf{M}')_{p}\) 同构（§2 第 7 节命题 18），我们的断言由下述引理得出：

引理 3. 设 B 是一个局部环，E 与 E' 是两个有限生成 B-模。若 E ≠ 0 且 E' ≠ 0，则 \(E \otimes_{B} E' \neq 0\)。

设 \(k\) 是 B 的剩余域。凭 §3 第 2 节命题 4，\(k \otimes_{\mathbf{B}} \mathbf{E} \neq 0\) 且 \(k \otimes_{\mathbf{B}} \mathbf{E}' \neq 0\)；于是我们推出

\[
(k \otimes_ {\mathbf {B}} \mathbf {E}) \otimes_ {k} (k \otimes_ {\mathbf {B}} \mathbf {E} ^ {\prime}) \neq 0
\]

（《代数》第 II 章 §3 第 7 节）。但由于张量积是结合的（同上，§3 第 8 节），该张量积同构于

\[
E \otimes_ {\mathbf {B}} ((k \otimes_ {k} k) \otimes_ {\mathbf {B}} \mathbf {E} ^ {\prime}) = \mathbf {E} \otimes_ {\mathbf {B}} (k \otimes_ {\mathbf {B}} \mathbf {E} ^ {\prime})
\]

从而同构于 \(k \otimes_{\mathbf{B}} (\mathbf{E} \otimes_{\mathbf{B}} \mathbf{E}')\)，由此得该引理。

推论. 设 M 是一个有限生成 A-模，n 是它的零化子。对 A 的每个理想 a，有 Supp(M/aM) = V(a) ∩ V(n) = V(a + n)。

\(\mathbf{M} / a\mathbf{M} = \mathbf{M}\otimes_{\mathbf{A}}(\mathbf{A} / a)\)，且 \(\mathbf{A} / a\) 是有限生成的。

命题 19. 设 A、B 是两个环，\(\phi: \mathbf{A} \to \mathbf{B}\) 是一个同态，且

\[
^ {a} \phi : \operatorname{Spec} (\mathbf {B}) \rightarrow \operatorname{Spec} (\mathbf {A})
\]

与 \(\phi\) 相伴的连续映射（命题 13）。对每个 A-模 M，\(\operatorname{Supp}(\mathbf{M}_{(\mathbf{B})}) \subset {}^a\overline{\phi}^{-1}(\operatorname{Supp}(\mathbf{M}))\)；若此外 M 还是有限生成的，则

\[
\operatorname{Supp} (\mathbf {M} _ {(\mathbf {B})}) = ^ {a} \bar {\phi} ^ {- 1} (\operatorname{Supp} (\mathbf {M})).
\]

设 \(q\) 是 \(B\) 的一个素理想，且 \(p = \overline{\phi}^{-1}(q)\)。设 \(q\) 属于 \(\operatorname{Supp}(M_{(B)})\)；则 \(M_{(B)} \otimes_B B_q = (M \otimes_A B) \otimes_B B_q = M \otimes_A B_q = (M \otimes_A A_p) \otimes_A B_q\)，因为同态 \(A \to B \to B_q\) 分解为 \(A \to A_p \to B_q\)（\(\S 2\) 第 1 节命题 2）；于是假设 \(M_{(B)} \otimes_B B_q \neq 0\) 蕴含 \(M \otimes_A A_p \neq 0\)，由此得第一个断言。由于同态 \(\phi_q: A_p \to B_q\) 是局部的，第二个断言由下述引理得出：

引理 4. 设 A、B 是两个局部环，\(\rho: \mathbf{A} \to \mathbf{B}\) 是一个局部同态，E 是一个有限生成 A-模。若 \(\mathbf{E} \neq 0\)，则 \(\mathrm{E}_{(\mathbf{B})} \neq 0\)。

设 m 是 A 的极大理想，k = A/m 是剩余域；假设蕴含 \(B \otimes_{A} k = B/mB \neq 0\)；由于张量积是结合的，\((\mathbf{E} \otimes_{\mathbf{A}} \mathbf{B}) \otimes_{\mathbf{A}} k\) 同构于 \(\mathbf{E} \otimes_{\mathbf{A}} (\mathbf{B} \otimes_{\mathbf{A}} k)\)，从而同构于 \(\mathbf{E} \otimes_{\mathbf{A}} (k \otimes_{k} (\mathbf{B} \otimes_{\mathbf{A}} k))\)，最终同构于 \((\mathbf{E} \otimes_{\mathbf{A}} k) \otimes_{k} (\mathbf{B} \otimes_{\mathbf{A}} k)\)；凭 §3 第 2 节命题 4，\(E \otimes_{A} k \neq 0\)，因而 \((\mathbf{E} \otimes_{\mathbf{A}} \mathbf{B}) \otimes_{\mathbf{A}} k \neq 0\)（《代数》第 II 章 §3 第 7 节），更加有 \(E \otimes_{A} B \neq 0\)。

命题 20. 设 A 是一个环，M 是一个有限生成 A-模。对每个素理想 \(\mathfrak{p} \in \operatorname{Supp}(\mathbf{M})\)，存在非零 A-同态 \(w: \mathbf{M} \to \mathbf{A}/\mathfrak{p}\)。

设 \(\mathfrak{p} \in \operatorname{Supp}(\mathbf{M})\)。由于 \(\mathbf{M}\) 是有限生成的且 \(\mathbf{M}_{\mathfrak{p}} \neq 0\)，

\[
\mathbf {M} _ {\mathfrak {p}} / \mathfrak {p} \mathbf {M} _ {\mathfrak {p}} = \mathbf {M} _ {\mathfrak {p}} \otimes_ {\mathbf {A}} (\mathbf {A} _ {\mathfrak {p}} / \mathfrak {p} \mathbf {A} _ {\mathfrak {p}}) \neq 0
\]

（§3 第 2 节命题 4）。设 \(\mathbf{K} = \mathbf{A}_{\mathfrak{p}} / \mathfrak{p}\mathbf{A}_{\mathfrak{p}}\) 是整环 \(\mathbf{A} / \mathfrak{p}\) 的分式域；由于 \(\mathbf{M}_{\mathfrak{p}} / \mathfrak{p}\mathbf{M}_{\mathfrak{p}}\) 是 \(\mathbf{K}\) 上的一个向量空间且不化为 0，存在非零线性型 \(u: \mathbf{M}_{\mathfrak{p}} / \mathfrak{p}\mathbf{M}_{\mathfrak{p}} \to \mathbf{K}\)。若 \((x_{i})_{1 \leqslant i \leqslant n}\) 是 \(\mathbf{M}, \bar{x}_{i}\) 的一个生成系（\(x_{i}\) 的像取在 \((\mathbf{A} / \mathfrak{p})\) -模 \(\mathbf{M}_{\mathfrak{p}} / \mathfrak{p}\mathbf{M}_{\mathfrak{p}}\) 中），则存在元素 \(\alpha \neq 0\)（属于 \(\mathbf{A} / \mathfrak{p}\)）使诸 \(\alpha u(\bar{x}_i)\) 属于 \(\mathbf{A} / \mathfrak{p}\)，对 \(1 \leqslant i \leqslant n\) 成立；因而 \(v = \alpha u\) 是非零的 \((\mathbf{A} / \mathfrak{p})\) -线性映射，从 \(\mathbf{M}_{\mathfrak{p}} / \mathfrak{p}\mathbf{M}_{\mathfrak{p}}\) 到 \(\mathbf{A} / \mathfrak{p}\)。复合

\[
w: \mathrm{M} \longrightarrow \mathrm{M} _ {p} \longrightarrow \mathrm{M} _ {p} / p \mathrm{M} _ {p} \stackrel {v} {\longrightarrow} \mathrm{A} / p
\]

因此就是所求的同态。

